%%%%%%
%
% $Autor: Wings $
% $Datum: 2019-03-05 08:03:15Z $
% $Pfad: Etcher $
% $Version: 4250 $
% !TeX spellcheck = en_GB
% !TeX encoding = utf8
% !TeX root = filename 
% !TeX TXS-program:bibliography = txs:///biber
%
%%%%%%


\chapter{Software Etcher}

\section{To Do}

\textcolor{red}{This chapter has to be reworked.}

\section{Introduction}


The software Etcher is Required for flashing the Jetson Nano system image onto the microSD card. Etcher, now known as balenaEtcher, is an open-source, cross-platform tool used for flashing disk images such as operating system images onto SD cards or USB drives. It's widely used to flash operating system images onto SD cards for devices like the Jetson Nano. \cite{Balena:2025}



\section{Key Features of Etcher}
\begin{itemize}
	\item \textbf{Cross-Platform Compatibility}: Etcher works on Windows, macOS, and Linux, making it accessible for developers using different systems.
	\item \textbf{Simple User Interface}: It offers a streamlined, easy-to-use interface for flashing images, making it ideal for beginners and advanced users alike.
	\item \textbf{Safe Flashing Process}: Etcher verifies the flashing process to ensure that the image is correctly written to the storage device, helping to prevent errors or corrupt files.
	\item \textbf{Support for Various Image Formats}: Etcher supports a wide range of image formats, including .img, .iso, and .zip files.
\end{itemize}

\section{How Etcher Works}

\begin{enumerate}
	\item \textbf{Select Image}: The user chooses the operating system image (e.g., Jetson Linux for Jetson Nano).
	\item \textbf{Select Target}: The user selects the SD card or USB drive where the image will be written.
	\item \textbf{Flashing Process}: Etcher writes the image to the storage device and verifies the process to ensure data integrity.
	\item \textbf{Validation}: Once the flashing process is complete, Etcher will perform a validation to ensure the image is correctly written without corruption.
\end{enumerate}

\section{Protocols Involved in Etcher}

\begin{itemize}
	\item \textbf{USB Protocol}: Etcher communicates with USB drives or SD cards using the USB protocol for data transfer.
	\item \textbf{File I/O Protocol}: Etcher uses standard file I/O protocols for reading and writing disk images to the selected device.
\end{itemize}

\subsection{Use Case}

\textbf{Flashing OS for Jetson Nano}: Etcher is used to flash the Jetson Linux operating system onto the Jetson Nano’s SD card to boot the device and set up the operating system. \cite{Balena:2025}



\section{Installing the Operating System}

To set up the Jetson Nano, a properly prepared microSD card is required. This involves downloading the operating system image file from the NVIDIA website and writing it to the microSD card using specialized software. The JetPack operating system is based on Linux and follows a Debian-like file system structure. Once the image is written to the microSD card, it can be inserted into the Jetson Nano for configuration. This section provides a step-by-step guide to installing the operating system.

\subsection{Writing the Operating System to the MicroSD Card}

The operating system image can be downloaded from the official NVIDIA website \cite{JetsonUserManual:2020, JetPackOS:2021}. The image is then written to the microSD card using tools like Etcher, available at \url{https://www.balena.io/etcher/} \cite{JetsonGettingStarted:2019}. 

\subsubsection{Formatting the MicroSD Card}

Before writing the image, the microSD card must be formatted. The procedure for formatting depends on the operating system of the host computer. Below are the steps for formatting the card on Windows:

\begin{enumerate}
	\item \textbf{Download and Install the SD Memory Card Formatter:} 
	Visit the SD Association's official website at \url{https://www.sdcard.org/downloads/formatter/} and download the SD Memory Card Formatter. Install and launch the application.
	
	\item \textbf{Format the MicroSD Card:}
	\begin{itemize}
		\item Insert the microSD card into your computer's SD card reader.
		\item Open the SD Memory Card Formatter. A dialogue window, as shown in Figure~\ref{JetPack:SdFormatierer}, will appear.
		\begin{figure}[h!]
			\centering
			\includegraphics{Etcher/WindowsSDCardFormatter}
			\caption{SD Memory Card Formatter Dialogue}\label{JetPack:SdFormatierer}
		\end{figure}
		\item Select the correct card drive.
		\item Choose the "Quick Format" option.
		\item Leave the "Volume Label" field unchanged.
		\item Click the "Format" button. A warning dialogue will appear. Confirm by clicking "Yes."
	\end{itemize}
\end{enumerate}

\subsubsection{Writing the Image File Using Etcher}

Once the microSD card is formatted, the JetPack image file can be written to it using Etcher:

\begin{enumerate}
	\item \textbf{Download and Launch Etcher:}
	Download Etcher from \url{https://www.balena.io/etcher/}. Install and launch the application. The main dialogue, shown in Figure~\ref{JetPack:Etcher}, will appear.
	\begin{figure}[h!]
		\centering
		\includegraphics[width=\textwidth]{Etcher/WindowsEtcher}
		\caption{Etcher Main Dialogue}\label{JetPack:Etcher}
	\end{figure}
	
	\item \textbf{Select the Image File:}
	Click the "Select image" button and choose the downloaded JetPack image file.
	
	\item \textbf{Handle Unformatted Cards:}
	If the microSD card is not formatted, Etcher will display an error, as shown in Figure~\ref{JetPack:NichtFormatiert}. Format the card before proceeding.
	\begin{figure}[h!]
		\centering
		\includegraphics{Etcher/WindowsEtcherCancel}
		\caption{Error Dialogue for Unformatted Card}\label{JetPack:NichtFormatiert}
	\end{figure}
	
	\item \textbf{Start Flashing:}
	Click the "Flash!" button to start writing the image file to the microSD card. The process may take around 10 minutes when using a USB 3.0 interface. After writing, Etcher will validate the image.
	
	\item \textbf{Ignore Windows Warnings:}
	After the process completes, Windows may display a message stating that the microSD card cannot be read, as shown in Figure~\ref{JetPack:NichtLesen}. This message can be ignored. Safely remove the card from the drive.
	\begin{figure}[h!]
		\centering
		\includegraphics{Etcher/WindowsSDCardPrompt}
		\caption{Windows Message for Unreadable MicroSD Card}\label{JetPack:NichtLesen}
	\end{figure}
\end{enumerate}

\subsubsection{Final Steps}

Once the microSD card is prepared, it can be inserted into the Jetson Nano. On the first boot, the system will guide you through initial configuration steps, including selecting a graphical or headless installation.

This process ensures a smooth setup and prepares the Jetson Nano for further development.

